% Bagian: model-theory
% Bab: interpolation
% Subbagian: interpolation-proof

\documentclass[../../../include/open-logic-section]{subfiles}

\begin{document}

\olfileid[id]{mod}{int}{prf}

\olsection{Teorema Interpolasi Craig}

\begin{thm}[Teorema Interpolasi Craig]
\ollabel{thm:interpol}
Jika $\Entails !A \lif !B$, maka terdapat !!a{sentence} $!C$ sedemikian
sehingga $\Entails !A \lif !C$ dan $\Entails !C \lif !B$, dan setiap
!!{constant}, !!{function}, serta !!{predicate} (selain $\eq$) dalam $!C$
terdapat dalam $!A$ maupun~$!B$. !!^{sentence} $!C$ disebut
\emph{interpolan} dari $!A$ dan~$!B$.
\end{thm}

\begin{proof}
Andaikan $\Lang{L}_1$ adalah bahasa dari $!A$ dan $\Lang{L}_2$ adalah
bahasa dari $!B$. Misalkan $\Lang{L}_0 = \Lang{L}_1 \cap \Lang{L}_2$.
Bagi setiap $i \in \{0, 1, 2 \}$, misalkan $\Lang{L}'_i$ diperoleh dari
$\Lang{L}_i$ dengan menambahkan tak hingga banyak !!{constant} baru
$\Obj c_0, \Obj c_1, \Obj c_2, \dots$.

Jika $!A$ tidak dapat dipenuhi, maka $\lexists[x][\eq/[x][x]]$ adalah suatu
interpolan. Jika $\lnot !B$ tidak dapat dipenuhi (sehingga $!B$ valid), maka
$\lexists[x][\eq[x][x]]$ adalah suatu interpolan. Karena itu, kita juga dapat
mengandaikan bahwa $!A$ dan $\lnot !B$ keduanya dapat dipenuhi.

Untuk membuktikan kontraposisi Teorema Interpolasi, andaikan tidak ada
interpolan bagi $!A$ dan $!B$. Dengan kata lain, andaikan $\{!A\}$ dan
$\{\lnot !B\}$ tidak terpisahkan dalam $\Lang{L}_0$.

Tujuan kita ialah memperluas pasangan $(\{ !A \}, \{\lnot!B\})$ menjadi
pasangan yang tidak terpisahkan secara maksimal $(\Gamma^*, \Delta^*)$.
Misalkan $!A_0$, $!A_1$, $!A_2$, \dots mencacah !!{sentence}s dari
$\Lang{L}'_1$, dan $!B_0$, $!B_1$, $!B_2$, \dots mencacah
!!{sentence}s dari~$\Lang{L}'_2$. Kita mendefinisikan dua barisan naik
himpunan !!{sentence}s $(\Gamma_n, \Delta_n)$, untuk $n \ge 0$, sebagai
berikut. Tetapkan $\Gamma_0 = \{ !A\}$ dan $\Delta_0 = \{\lnot !B \}$.
Dengan mengandaikan $(\Gamma_n, \Delta_n)$ telah didefinisikan, definisikan
$\Gamma_{n+1}$ dan $\Delta_{n+1}$ sebagai berikut:
\begin{enumerate}
\item Jika $\Gamma_n \cup \{!A_n \}$ dan $\Delta_n$ tidak terpisahkan
  dalam $\Lang{L}'_0$, masukkan $!A_n$ ke dalam $\Gamma_{n+1}$; jika
  keduanya terpisahkan, tetapkan $\Gamma_{n+1}=\Gamma_n$. Dalam kasus pertama
  tersebut, jika
  $!A_n$ adalah !!{formula} eksistensial $\lexists[x][!S]$, pilih suatu
  !!{constant} baru $c$ yang tidak terdapat dalam $\Gamma_n$, $\Delta_n$,
  $!A_n$, ataupun $!B_n$, lalu masukkan $\Subst{!S}{c}{x}$ ke dalam
  $\Gamma_{n+1}$.
\item Jika $\Gamma_{n+1}$ dan $\Delta_n \cup \{!B_n \}$ tidak
  terpisahkan dalam $\Lang{L}'_0$, masukkan $!B_n$ ke dalam
  $\Delta_{n+1}$; jika keduanya terpisahkan, tetapkan
  $\Delta_{n+1}=\Delta_n$. Dalam kasus pertama tersebut, jika $!B_n$ adalah
  !!{formula}
  eksistensial $\lexists[x][!S]$, pilih suatu !!{constant} baru $c$ yang
  tidak terdapat dalam $\Gamma_{n+1}$, $\Delta_n$, $!A_n$, ataupun $!B_n$,
  lalu masukkan $\Subst{!S}{c}{x}$ ke dalam $\Delta_{n+1}$.
\end{enumerate}
Terakhir, definisikan:
\begin{align*}
  \Gamma^* & = \bigcup_{n\ge 0} \Gamma_n, &
  \Delta^* & = \bigcup_{n\ge 0} \Delta_n.
\end{align*}
Dengan induksi serentak pada $n$, kini kita dapat membuktikan:
\begin{enumerate}
\item\ollabel{part-a} $\Gamma_n$ dan $\Delta_n$ tidak terpisahkan dalam
  $\Lang{L}'_0$;
\item\ollabel{part-b} $\Gamma_{n+1}$ dan $\Delta_n$ tidak terpisahkan
  dalam $\Lang{L}'_0$.
\end{enumerate}
Dasar bagi \olref{part-a} diberikan oleh \olref[sep]{lem:sep1}. Untuk bagian
\olref{part-b}, kita perlu membedakan tiga kasus:
\begin{enumerate}
\item Jika $\Gamma_0 \cup \{!A_0 \}$ dan $\Delta_0$ terpisahkan, maka
  $\Gamma_1 = \Gamma_0$ dan \olref{part-b} hanyalah \olref{part-a};
\item Jika $\Gamma_1 = \Gamma_0 \cup\{ !A_0\}$, maka $\Gamma_1$ dan
  $\Delta_0$ tidak terpisahkan berdasarkan konstruksi.
\item Tersisa kasus ketika $!A_0$ eksistensial, sehingga $\Gamma_1 =
  \Gamma_0 \cup \{ \lexists[x][!S], \Subst{!S}{c}{x} \}$. Berdasarkan
  konstruksi, $\Gamma_0 \cup \{ \lexists[x][!S]\}$ dan $\Delta_0$ tidak
  terpisahkan; maka, menurut \olref[sep]{lem:sep2}, $\Gamma_0 \cup \{
  \lexists[x][!S], \Subst{!S}{c}{x} \}$ dan $\Delta_0$ juga tidak
  terpisahkan.
\end{enumerate}
Ini menyelesaikan dasar induksi bagi \olref{part-a} dan \olref{part-b} di
atas. Sekarang tinjau langkah induksi. Untuk \olref{part-a}, jika
$\Delta_{n+1} = \Delta_n \cup \{ !B_n \}$, maka $\Gamma_{n+1}$ dan
$\Delta_{n+1}$ tidak terpisahkan berdasarkan konstruksi (bahkan ketika
$!B_n$ eksistensial, menurut \olref[sep]{lem:sep2}); jika $\Delta_{n+1} =
\Delta_n$ (karena $\Gamma_{n+1}$ dan $\Delta_n \cup \{!B_n\}$
terpisahkan), kita menggunakan hipotesis induksi pada \olref{part-b}. Untuk
langkah induksi bagi \olref{part-b}, jika $\Gamma_{n+2} = \Gamma_{n+1}
\cup \{!A_{n+1} \}$, maka $\Gamma_{n+2}$ dan $\Delta_{n+1}$ tidak
terpisahkan berdasarkan konstruksi (bahkan ketika $!A_{n+1}$ eksistensial,
menurut \olref[sep]{lem:sep2}); jika $\Gamma_{n+2} = \Gamma_{n+1}$, kita
menggunakan kasus induksi bagi \olref{part-a} yang baru saja dibuktikan. Ini
menyelesaikan induksi pada \olref{part-a} dan \olref{part-b}.

Akibatnya, $\Gamma^*$ dan $\Delta^*$ tidak terpisahkan. Jika tidak, menurut
kekompakan terdapat suatu kalimat pemisah dan suatu $n \ge 0$ sehingga
kalimat itu memisahkan $\Gamma_n$ dan $\Delta_n$, bertentangan dengan
\olref{part-a}. Khususnya, $\Gamma^*$ dan $\Delta^*$ konsisten: sebab jika
yang pertama atau yang kedua tidak konsisten, keduanya masing-masing akan
dipisahkan oleh $\lexists[x][\eq/[x][x]]$ atau
$\lforall[x][\eq[x][x]]$.

Sekarang kita tunjukkan bahwa $\Gamma^*$ konsisten maksimal dalam
$\Lang{L}'_1$, dan demikian pula $\Delta^*$ dalam $\Lang{L}'_2$. Untuk
yang pertama, andaikan $!A_n \notin \Gamma^*$ dan $\lnot !A_n \notin
\Gamma^*$ bagi suatu $n \ge 0$. Jika $!A_n \notin \Gamma^*$, maka
$\Gamma_n \cup \{!A_n \}$ terpisahkan dari $\Delta_n$, sehingga terdapat
$!C \in \Lang{L}'_0$ sedemikian sehingga keduanya berlaku:
\begin{align*}
  \Gamma^* & \Entails !A_n \lif !C, &
  \Delta^* & \Entails \lnot !C.
\end{align*}
Demikian pula, jika $\lnot !A_n \notin \Gamma^*$, terdapat $!C' \in
\Lang{L}'_0$ sedemikian sehingga keduanya berlaku:
\begin{align*}
  \Gamma^* & \Entails \lnot !A_n \lif !C', &
  \Delta^* & \Entails \lnot !C'.
\end{align*}
Menurut logika proposisional, $\Gamma^* \Entails !C \lor !C'$ dan
$\Delta^* \Entails \lnot (!C \lor !C')$, sehingga $!C \lor !C'$
memisahkan $\Gamma^*$ dan $\Delta^*$. Argumen serupa menetapkan bahwa
$\Delta^*$ maksimal.

Terakhir, kita tunjukkan bahwa $\Gamma^* \cap \Delta^*$ konsisten maksimal
dalam $\Lang{L}'_0$. Himpunan ini jelas konsisten karena merupakan irisan
himpunan-himpunan konsisten. Untuk menunjukkan kemaksimalannya, misalkan
$!S \in \Lang{L}'_0$. Sekarang, $\Gamma^*$ maksimal dalam $\Lang{L}'_1
\supseteq \Lang{L}'_0$, dan demikian pula $\Delta^*$ maksimal dalam
$\Lang{L}'_2 \supseteq \Lang{L}'_0$. Akibatnya, $!S \in \Gamma^*$ atau
$\lnot !S \in \Gamma^*$, dan $!S \in \Delta^*$ atau $\lnot !S \in
\Delta^*$. Jika $!S \in \Gamma^*$ dan $\lnot !S \in \Delta^*$, maka
$!S$ akan memisahkan $\Gamma^*$ dan $\Delta^*$; dan jika $\lnot !S \in
\Gamma^*$ dan $!S \in \Delta^*$, maka $\Gamma^*$ dan $\Delta^*$ akan
dipisahkan oleh $\lnot !S$. Karena itu, $!S \in \Gamma^* \cap \Delta^*$
atau $\lnot !S \in \Gamma^* \cap \Delta^*$, dan $\Gamma^* \cap
\Delta^*$ maksimal.

Karena $\Gamma^*$ konsisten maksimal, himpunan ini mempunyai model
$\Struct{M}'_1$ yang !!{domain}nya, $\Domain{M'_1}$, terdiri atas tepat
semua unsur $\Assign{c}{M'_1}$ yang menafsirkan !!{constant}s---seperti
dalam bukti teorema kelengkapan
(\olref[fol][com][cth]{thm:completeness}). Demikian pula, $\Delta^*$
mempunyai model $\Struct{M}'_2$ yang !!{domain}nya, $\Domain{M'_2}$,
diberikan oleh interpretasi $\Assign{c}{M'_2}$ dari !!{constant}s.

Misalkan $\Struct{M_1}$ diperoleh dari $\Struct{M'_1}$ dengan membuang
interpretasi bagi !!{constant}s, !!{function}s, dan !!{predicate}s dalam
$\Lang{L}'_1 \setminus \Lang{L}'_0$, dan lakukan hal serupa untuk
$\Struct{M_2}$. Maka pemetaan $h \colon M_1 \to M_2$ yang didefinisikan
oleh $h(\Assign{c}{M'_1}) = \Assign{c}{M'_2}$ merupakan isomorfisme dalam
$\Lang{L}'_0$, karena $\Gamma^* \cap \Delta^*$ konsisten maksimal dalam
$\Lang{L}'_0$, sebagaimana telah ditunjukkan. Hal ini mengikuti karena setiap
$\Lang{L}'_0$-!!{sentence} termasuk dalam $\Gamma^*$ dan $\Delta^*$
sekaligus, atau tidak termasuk dalam keduanya: sehingga
$\Assign{c}{M'_1} \in \Assign{P}{M'_1}$ jika dan hanya jika
$\Atom{P}{c} \in \Gamma^*$ jika dan hanya jika $\Atom{P}{c} \in
\Delta^*$ jika dan hanya jika $\Assign{c}{M'_2} \in \Assign{P}{M'_2}$.
Syarat-syarat isomorfisme lainnya dapat ditetapkan dengan cara serupa.

Sekarang kita definisikan suatu model $\Struct{M}$ bagi !!{language}
$\Lang{L}'_1 \cup \Lang{L}'_2$ sebagai berikut:
\begin{enumerate}
\item !!^{domain} $\Domain{M}$ tidak lain adalah $\Domain{M_2}$, yakni
  himpunan semua unsur $\Assign{c}{M'_2}$;
\item Jika !!a{predicate}~$P$ berada dalam $\Lang{L}'_2 \setminus
  \Lang{L}'_1$, maka $\Assign{P}{M} = \Assign{P}{M'_2}$;
\item Jika suatu predikat $P$ berada dalam $\Lang{L}'_1\setminus
  \Lang{L}'_2$, maka $\Assign{P}{M} = h(\Assign{P}{M'_1})$, yakni
  $\tuple{\Assign{c_1}{M'_2}, \dots, \Assign{c_n}{M'_2}} \in
  \Assign{P}{M}$ jika dan hanya jika $\tuple{\Assign{c_1}{M'_1}, \dots,
  \Assign{c_n}{M'_1}} \in \Assign{P}{M'_1}$.
\item Jika !!a{predicate} $P$ berada dalam $\Lang{L}'_0$, maka
  $\Assign{P}{M} = \Assign{P}{M'_2} = h(\Assign{P}{M'_1})$.
\item !!^{function}s dari $\Lang{L}'_1 \cup \Lang{L}'_2$, termasuk
  !!{constant}s, ditangani dengan cara serupa.
\end{enumerate}

Terakhir, dengan induksi pada !!{formula}s ditunjukkan bahwa $\Struct{M}$
bersepakat dengan $\Struct{M'_1}$ pada semua !!{formula}s dari
$\Lang{L}'_1$ dan dengan $\Struct{M'_2}$ pada semua !!{formula}s dari
$\Lang{L}'_2$. Khususnya, $\Struct{M} \models \Gamma^* \cup \Delta^*$,
sehingga $\Struct{M} \models !A$ dan $\Struct{M} \models \lnot!B$, dan
$\not\Entails !A \lif !B$. Ini menyelesaikan bukti Teorema Interpolasi
Craig.
\end{proof}

\end{document}
